\documentclass[10pt,a4paper]{amsart}
\usepackage[margin=19mm]{geometry}
\usepackage{amsmath,amssymb,mathtools}
\usepackage[scr=rsfs,cal=euler]{mathalfa}
\usepackage{xfrac}
\usepackage[unicode,hidelinks,bookmarks=false]{hyperref}
\newcommand{\ie}{i.e.\ }
\newcommand{\cf}{cf.\ }
\newcommand{\resp}{respectively\ }
\newcommand{\Subsection}{\subsection}
\providecommand{\nobreakdash}{\nobreak\hbox{-}}
\setlength{\emergencystretch}{2em}
\multlinegap0pt
\usepackage[shortlabels]{enumitem}
\usepackage{amssymb}
\usepackage{tikz}
\usepackage{dsfont}
\usepackage{float}
\newtheorem{lem}{Lemma}
\newcommand{\De}{\Delta}
\newcommand{\ba}{\begin{equation} \begin{aligned}}
\newcommand{\bs}[1]{\boldsymbol{#1}}
\newcommand{\ea}{\end{aligned} \end{equation}}
\newcommand{\rhob}{{\bar{\rho}}}
\title{Maximally heavy dynamics in the causal diamond}
\author{David Poland, Gordon Rogelberg}
\date{Source: arXiv:2603.28853v1 (CC BY 4.0)}
\newcommand{\cO}{\mathcal O}
\newcommand{\bh}{{\bar{h}}}
\begin{document}
\maketitle
\par\medskip\noindent\textbf{Current collection classification: H1.} Retained author material; this item does not qualify for the advanced H2/H3 collection. See the canonical metadata for the source-specific rationale.\par\medskip


{\sloppy \noindent\emph{Source and licence.} Maximally heavy dynamics in the causal diamond. Version arXiv:2603.28853v1 (CC BY 4.0), distributed by arXiv under the Creative Commons Attribution 4.0 International licence, \url{http://creativecommons.org/licenses/by/4.0/}, as recorded in the arXiv licence metadata for this exact version and shown on the abstract page \url{https://arxiv.org/abs/2603.28853v1.} Full licence notice: \path{licenses/arxiv-2603.28853-CC-BY-4.0.txt}.\par\smallskip\par}

\noindent\textit{Editorial scope.} Counted unit: the article's lem sharpmomentbound (source0.tex lines 1895--1901). The article's complete original proof of it is delivered with it (source0.tex lines 1902--1923, one \texttt{proof} environment of 22 lines). Additional author definitions and conventions recorded in the correction metadata are retained below as uncounted context. Source caveats and pending qualification issues are recorded separately. Author statements and proofs remain unchanged; the fixed complete author source remains bundled.\par\medskip
\par\medskip\noindent\textit{Retained author context (same fixed source; added for the selected problem).}\par
In this section, we review the basic setup of our problem and underline its parallel to the torus partition function in 2d CFT. Since this is mostly review, expert readers may skip this section. Consider a correlation function of four identical scalar operators $\cO$ with scaling dimension $\De$ in a unitary CFT on $\widehat{\mathbb{R}}^d =  \mathbb{R}^d\bigcup\{\infty\} \cong S^d $:
\ba
\langle \cO(x_1) \cO(x_2) \cO(x_3) \cO(x_4)\rangle = \langle \cO(x_1) \cO(x_2) \rangle \langle \cO(x_3) \cO(x_4) \rangle \mathcal{G}(u,v),
\label{cft4ptfunction}
\ea
where $\mathcal{G}(u,v)$ denotes the reduced correlator, which is a function of only cross ratios. Since $u,v$ are symmetric polynomials in $z,\bar{z}$, reduced correlators have a chiral symmetry of $\mathcal{G}(\bs{z}) = \mathcal{G}(\bar{\bs{z}})$. For all $x^\mu_i \in \widehat{\mathbb{R}}^d$, $u,v$ are positive and $\bar{z} = z^*$.

Factoring out the disconnected kinematic prefactor, the reduced correlator admits the following $\bs{h} \leftrightarrow\bar{\bs{h}}$ symmetric decomposition into radial monomials:
\ba
\mathcal{G}(\bs{\rho}) = \sum_{\bs{h}} a_{\bs{h}} \bs{\rho}^{\bs{h}} = \sum_{\bs{h}} a_{\bar{\bs{h}}} \bs{\rho}^{\bs{h}},
\label{monomialdecomposition}
\ea
where $|\lambda_{h,\bh}|^2 16^{-\De} = a_{\bs{h}}>0$ are real-positive with $a_{\bs{0}} = 1$, and the second equality follows from chiral symmetry. More geometrically, chiral symmetry arises from the equivalence of radially quantizing around the origin or the point at infinity in $\widehat{\mathbb{R}}^d$, at the cost of flipping $J \to -J$, resulting in the exchange of $H \leftrightarrow \bar{H}$.
 
The associativity of the OPE implies that $\mathcal{G}$ is invariant under permutations of the coordinates $\{x_i\}$, giving rise to non-trivial sum rules (crossing equations) relating different OPE decompositions. Of these, only the s-t crossing equation provides usable constraints within the causal diamond, since the u-channel decomposition does not converge in $\mathscr{C}$ due to the absence of co-dimension 1 spheres separately enclosing the operator pairs at $(x_1,x_3)$ and $(x_2,x_4)$. In this work, we therefore only impose s-t crossing symmetry, given by the following constraint:
\ba
\label{stcrossing}
\mathcal{G}(\bs{\rho}) = \left|\frac{4 \bs{\rho}}{(1-\bs{\rho})^2}\right|^{2\De}\mathcal{G}(\hat{\bs{\rho}}),
\ea
where $\hat{\bs{\rho}} = \left(\frac{1-\sqrt{\bs{\rho}}}{1+\sqrt{\bs{\rho}}} \right)^2$ is the image of $\bs{\rho}$ under $\bs{z} \to 1-\bs{z}$.

We now review the basic formalism for recasting the OPE decomposition as an integral over a positive definite measure. Consider the radial monomial decomposition in (\ref{monomialdecomposition}) and fix some kinematics in the causal diamond $\boldsymbol{\rho} = (\rho,\bar{\rho})\in(0,1)^2$. We define the normalized OPE measure at point $\bs{\rho} \in \mathscr{C}$ as
\ba
d\mu(\bs{\rho}; \bs{h}) = \frac{1}{\mathcal{G}(\bs{\rho})}\sum_{\bs{h}'} a_{\bs{h}'} \bs{\rho}^{\bs{h}'} \delta^{(2)}(\bs{h} - \bs{h}') d\bs{h}.
\ea
We can now rewrite the radial monomial decomposition as 
\ba
\frac{\mathcal{G}\left(\bs{\rho}\right)}{\mathcal{G}\left(\bs{\rho}_0\right)} = \int_{\mathbb{R}^2_+} \left(\frac{\bs{\rho}}{\bs{\rho}_0}\right)^{\bs{h}} d\mu(\bs{\rho}_0 ; \bs{h}) .
\label{radmonomialdecomp}
\ea
This integral reproduces exactly the original sum and therefore absolutely converges in the unit disk of $\rho,\bar{\rho}$.

We denote the normalized moments of the OPE measure at base point $\bs{\rho}$ as
\ba
\omega_{j,k}(\bs{\rho}) = \int h^j \bh^k d\mu(\bs{h};\bs{\rho}).
\ea

\begin{lem}[Sharp bound on the first moment in the causal diamond]
\label{sharpmomentbound}
The first normalized moments satisfy the sharp bounds
\ba
\omega_{1,0}(\bs{\rho}) \leq \De \frac{1+\rho}{1-\rho} \quad \text{and} \quad \omega_{0,1}(\bs{\rho}) \leq \De \frac{1+\rhob}{1-\rhob} \quad \forall \bs{\rho} \in (0,1)^2.
\ea
\end{lem}

\begin{proof}
The moments satisfy $\omega_{j,k}(\bs{\rho}) = \omega_{k,j}(\bar{\bs{\rho}})$, so it suffices to check for $(j,k) = (1,0)$ as long as the bound holds for all $\bs{\rho} \in (0,1)^2$. We write $\bs{\rho} = (e^{-\bs{s}}) = (e^{-s},e^{-\bar{s}})$ and take the logarithm of the crossing equation
\ba
\log(\mathcal{G}(\bs{s})) = \De \log(|\bs{\chi}(\bs{s})|) + \log( \mathcal{G}( \hat{\bs{s}} )),
\ea
where 
\ba
\hat{\bs{s}} = \left(- \log(\hat{\bs{\rho}}(\bs{s}))\right) = \left(- 2 \log\left(\frac{1-\sqrt{e^{-s}}}{1+\sqrt{e^{-s}}}\right) ,- 2 \log\left(\frac{1-\sqrt{e^{-\bar{s}}}}{1+\sqrt{e^{-\bar{s}}}}\right) \right)
\ea
and 
\ba
|\bs{\chi}(\bs{s})| = \left(  \frac{16 e^{-(s+\bar{s})}}{(1-e^{-s})^2 (1-e^{-\bar{s}})^2}\right).
\ea
We now take a derivative of both sides with respect to $-s$ and evaluate at $\bs{s} = \left(-\log(\bs{\rho}) \right)$ to compute the first moment $\omega_{1,0}(\bs{\rho})$ on the LHS
\ba
\omega_{1,0}(\bs{\rho}) = \De \partial_{-s} \log(|\bs{\chi}(\bs{s})|)|_{\bs{s} = -\log(\bs{\rho})} + \left(\partial_{-s} \hat{\rho}(s) \right)|_{s = -\log(\rho)} \omega_{1,0}(\hat{\bs{\rho}}).
\ea
By the positivity of the support of the OPE measure, $\omega_{1,0}(\bs{\rho})$ is positive for all points in the causal diamond, as is $\omega_{1,0}(\hat{\bs{\rho}})$. It is also easy to check that $\partial_{-s} \hat{\rho}(s) = \frac{2 \sqrt{e^{-s}}}{e^{-s} -1} \leq 0$ for all points in the causal diamond, so the second term on the RHS is always negative. Thus, we have
\ba
\omega_{1,0}(\bs{\rho}) &= \De \partial_{-s} \log(|\bs{\chi}(\bs{s})|)|_{\bs{s} = -\log(\bs{\rho})} + \left(\partial_{-s} \hat{\rho}(s) \right)|_{s = -\log(\rho)} \omega_{1,0}(\hat{\bs{\rho}})\\& \leq \De \partial_{-s} \log(|\bs{\chi}(\bs{s})|)|_{\bs{s} = -\log(\bs{\rho})}  = \De \frac{1+\rho}{1-\rho}.
\ea
\end{proof}

\end{document}
